\documentclass[11pt,a4paper]{article}
\usepackage[margin=2.3cm]{geometry}
\usepackage{amsmath,amssymb,graphicx,booktabs,xcolor,caption,microtype}
\usepackage[colorlinks=true,linkcolor=NavyBlue,citecolor=Green,urlcolor=NavyBlue]{hyperref}
\usepackage[usenames,dvipsnames]{xcolor}
\usepackage{tcolorbox}
\usepackage{enumitem}
\tcbuselibrary{skins,breakable}

\newtcolorbox{keybox}[1][]{colback=Goldenrod!10,colframe=Goldenrod!70!black,
  boxrule=0.7pt,left=6pt,right=6pt,top=5pt,bottom=5pt,breakable,#1}
\newtcolorbox{warnbox}[1][]{colback=BrickRed!6,colframe=BrickRed!70!black,
  boxrule=0.7pt,left=6pt,right=6pt,top=5pt,bottom=5pt,breakable,#1}

\newcommand{\Sop}{S_{\mathrm{op}}}
\newcommand{\chreq}{\chi_{\mathrm{req}}}
\newcommand{\nsat}{n_{\mathrm{sat}}}

\title{\bfseries Where does the dissipative Heisenberg chain become\\
classically hard?\\[4pt]
\large Operator entanglement and bond-dimension requirements for\\
vectorised MPS/TEBD, and what they imply for hardware runs}
\author{Internal working note}
\date{\today}

\begin{document}
\maketitle

\begin{abstract}
\noindent
We map out the region of parameter space $(n,\gamma,t)$ in which a vectorised
matrix-product simulation of the disordered Heisenberg chain under bulk
amplitude damping stops being cheap, with the goal of choosing parameters for a
quantum-hardware experiment that implements the damping via ancillas. Three
results. (i)~Amplitude damping produces an \emph{entanglement barrier}: the
operator entanglement rises, peaks, and decays to zero, because the dissipator's
steady state is the trivial product state $|0\cdots0\rangle$. The peak sits at
$t^\ast\!\approx\!0.45\,n$, \emph{not} at $t\!\sim\!1/\gamma$ as we initially
assumed. (ii)~Above a threshold system size $\nsat(\gamma)\approx1.4/\gamma$ the
peak becomes completely independent of $n$ --- an area law --- and the classical
cost stops growing, making the experiment pointless at any size. This is
measured, not inferred: at $\gamma=0.10$ the peak is $4.86,4.89,4.90,4.90$~bits
for $n=12,16,20,24$. (iii)~Along the operating line $\gamma^\ast=1.4/n$,
$t^\ast=0.45n$, the required bond dimension already exceeds single-node
feasibility at $n\approx20$. The useful empirical law is
$\chreq \approx 6\cdot 2^{1.61\,\Sop}$: the Schmidt tail costs a further
factor $2^{0.6\,\Sop}$ beyond the naive $2^{\Sop}$.
We also report a methodological finding: the order-2 product formula in the
existing codebase is not second order in the open case, and we give a corrected
symmetric splitting. We close by stating plainly what these results do
\emph{not} establish --- they bound one classical method, and the obvious
competitor (quantum trajectories) is untested.
\end{abstract}

\tableofcontents
\vspace{4mm}

%======================================================================
\section{Purpose and scope}

The target experiment is a quantum circuit that implements the Heisenberg
Hamiltonian together with an amplitude-damping channel realised through ancilla
qubits. For that experiment to be worth running, the parameters
$(n,\gamma,t)$ must be chosen in a region where a classical simulation is
\emph{not} cheap. This note determines that region.

The scope is deliberately narrow. We study only the physical state
$\rho(t)$ and the cost of representing it as a matrix-product operator. Nothing
here concerns the middle-out contraction scheme, the $F_{ij}$ Gram matrix, or
multi-product formulas; those objects have a different and much less favourable
cost scaling, which is why the state route reaches $n=24$ where the $F$ route
cannot.

The reader who wants only the answer should read
Sec.~\ref{sec:operating} (recommended operating points) and
Sec.~\ref{sec:limits} (what this does not prove).

%======================================================================
\section{The system}

\subsection{Hamiltonian and dissipator}

We simulate the Lindblad master equation
\begin{equation}
  \frac{d\rho}{dt} \;=\; -i[H,\rho]
  \;+\; \sum_{j=1}^{n}\gamma_j\Big(\sigma^-_j\rho\,\sigma^+_j
        -\tfrac12\{\sigma^+_j\sigma^-_j,\rho\}\Big),
  \label{eq:lindblad}
\end{equation}
with the disordered anisotropic Heisenberg Hamiltonian of
Eq.~(13) of the main manuscript,
\begin{equation}
  H \;=\; -\sum_{j=1}^{n-1}\Big[
      J_j\big(X_jX_{j+1}+Y_jY_{j+1}\big) + 2J_j\,Z_jZ_{j+1}\Big].
  \label{eq:H}
\end{equation}
The dissipator contains \emph{only} the lowering operator $\sigma^-$: this is
amplitude damping, not a balanced emission/absorption bath. Two consequences
follow immediately and shape everything below.

\begin{itemize}[leftmargin=1.4em,itemsep=2pt]
\item \textbf{The steady state is trivial.} Eq.~\eqref{eq:lindblad} drives the
  chain to $\rho_\infty=|0\cdots0\rangle\langle0\cdots0|$, a product state with
  zero entanglement of any kind. Whatever entanglement the dynamics generates
  must therefore come back down. There is no long-time plateau to exploit, and
  ``run longer and it gets harder'' is false for this model.
\item \textbf{The noise is non-unital.} Intuition imported from depolarising
  noise does not transfer; see Ref.~\cite{fefferman2024nonunital}, which shows
  that under non-unital noise the output distribution of a random circuit never
  anticoncentrates at any depth, in contrast to the unital case.
\end{itemize}

Unless stated otherwise all runs use a uniform chain $J_j=J=0.5$ (the mean of
the $\mathcal{U}[1/4,3/4]$ disorder used elsewhere in the project) and uniform
$\gamma_j=\gamma$, starting from the N\'eel product state. A uniform $J$ was
chosen deliberately for the scaling study: a single disorder realisation adds
run-to-run noise that is easily mistaken for $n$-dependence.

\subsection{What the gates actually implement}
\label{sec:gateconv}

A caveat that matters for anyone converting our times into the units of the
literature. The two-site bond gate in the codebase is
$R_{ZZ}(\beta)R_{YY}(\alpha)R_{XX}(\alpha)$ with $\alpha=J_j\,dt$ and
$\beta=2J_j\,dt$, which is $\exp(-i\,dt\,h_j)$ for
\begin{equation}
  h_j \;=\; \tfrac{J_j}{2}\big(X_jX_{j+1}+Y_jY_{j+1}\big) + J_j\,Z_jZ_{j+1}
  \;=\; -\tfrac12 H_j^{\text{Eq.}\eqref{eq:H}} .
\end{equation}
The overall sign is irrelevant for entanglement, but the factor of two is not:
the effective time is half what a naive reading of Eq.~\eqref{eq:H} suggests.
All times quoted in this note are the code's $t$. The anisotropy ratio is
$|J_z/J_{xy}|=2$, which is on the mild side --- Ref.~\cite{preisser2023comparing}
finds that operator entanglement \emph{decreases} with increasing $ZZ$ strength.

%======================================================================
\section{Simulation method}

\subsection{Vectorisation}

We vectorise $\rho\to|\rho\rangle\rangle$ and propagate it as an MPS over $n$
doubled sites, each carrying a ket and a bra index, so the local dimension is
$4$. Time evolution is TEBD: alternating odd/even Hamiltonian layers plus a
single-site dissipator layer, with singular-value truncation at every two-site
gate.

\subsection{Bond-dimension ceilings}

Three different ceilings appear in this project and confusing them is easy:
\begin{center}
\begin{tabular}{lll}
\toprule
Object & Local dim. & Ceiling at bond $b$ \\
\midrule
$|\psi\rangle$, pure state MPS & 2 & $2^{\min(b,\,n-b)}$ \\
$|\rho\rangle\rangle$, MPDO as MPS & 4 & $4^{\min(b,\,n-b)}$ \\
$F_{ij}$, superoperator MPO & 16 & $16^{\min(b,\,n-b)}$ \\
\bottomrule
\end{tabular}
\end{center}
At $n=20$ these are $10^3$, $10^6$ and $10^{12}$. The state route is what makes
this study possible at all.

\subsection{Trotter splitting, and a correction}
\label{sec:trotter}

The existing \texttt{get\_open\_step\_gates} at \texttt{order=2} composes, in
application order,
\begin{equation*}
  \mathrm{odd}(dt/2),\quad \mathrm{even}(dt),\quad \mathrm{diss}(dt),\quad
  \mathrm{odd}(dt/2)
  \qquad\text{i.e.}\qquad e^{A\,dt/2}e^{C\,dt}e^{B\,dt}e^{A\,dt/2},
\end{equation*}
with $A,B,C$ the odd, even and dissipator generators. The outer layer is symmetrised but $B$ and $C$ are
composed as a bare product, leaving an uncancelled commutator,
\begin{equation}
  e^{Cdt}e^{Bdt} = \exp\!\Big((B+C)dt + \tfrac12[C,B]\,dt^2+\dots\Big),
\end{equation}
so the global error carries an $O(dt)$ term whenever $[B,C]\neq0$. For
$\gamma=0$ we have $C=0$, the composition collapses to plain Strang and is
genuinely second order --- which is why closed-system testing would never have
revealed this.

We measured the order directly at $n=8$, where the MPDO ceiling $4^4=256$ equals
the truncation cap and the run is therefore \emph{exact}, so the only error left
is the product formula. Using ratios of successive differences (which need no
converged reference), Table~\ref{tab:order} and Fig.~\ref{fig:method}(b):

\begin{table}[h]
\centering
\caption{Effective Trotter order $p_{\mathrm{eff}}=\log_2$(ratio of successive
differences on halving $dt$), measured at $n=8$, $\gamma=0.05$, with truncation
identically zero. A clean second-order scheme gives $p_{\mathrm{eff}}=2$.}
\label{tab:order}
\begin{tabular}{lcccc}
\toprule
$t$ & 2 & 4 & 6 & 10 \\
\midrule
\texttt{:project} (existing) & 1.61,\,1.44 & 1.57,\,1.39 & 2.30,\,2.94 & sign change \\
\texttt{:strang} (corrected) & 2.02,\,2.20 & 1.97,\,1.96 & 2.00,\,2.04 & 2.09,\,2.53 \\
\bottomrule
\end{tabular}
\end{table}

The existing scheme is not cleanly first order either; it is a mixture, with the
error changing \emph{sign} between $dt=0.05$ and $dt=0.025$ at $t=10$ --- the
fingerprint of two error terms of opposite parity. The corrected palindromic
splitting
\begin{equation*}
  \mathrm{odd}(dt/2),\ \mathrm{even}(dt/2),\ \mathrm{diss}(dt),\
  \mathrm{even}(dt/2),\ \mathrm{odd}(dt/2)
\end{equation*}
is cleanly second order at every time. Being palindromic ($S(-dt)=S(dt)^{-1}$)
forces the effective generator to expand in even powers of $dt$ only.

Practically: at $t=4$, \texttt{:strang} at $dt=0.05$ is more accurate than
\texttt{:project} at $dt=0.025$ while taking a quarter of the steps. It costs
$\sim$33\% more per step (one extra even layer), a net factor $\sim$2.6 at fixed
accuracy. All production runs reported below use \texttt{:strang} at $dt=0.05$,
for which the error in $\Sop$ relative to $dt=0.02$ is
0.02\%/0.16\%/0.50\% at $t=3/6/9$.

\begin{warnbox}
\textbf{Wider implication, not pursued here.} The order-4 routine applies the
Yoshida triple-jump to this same order-2 composition. That construction reaches
fourth order only if the base is second order \emph{and} palindromic, and the
base is neither. Anything in the project that relies on the nominal order of
\texttt{get\_open\_step\_gates} --- in particular multi-product-formula
coefficient derivations, which assume an even-powers-only error expansion ---
should be re-examined.
\end{warnbox}

\subsection{The closed-system baseline}

The $\gamma=0$ control must be computed as a \emph{pure state}, not as a density
matrix. Our first attempt propagated $|\rho\rangle\rangle$ with
\texttt{dissipation=false} at $\chi_{\max}=256$ and produced purities
$\mathrm{Tr}\rho^2 = 1.69,\ 4.19,\ 24.2$ at $n=12,16,20$. Purity is identically
$1$ for a pure state, so those runs were unphysical beyond $t\approx2.5$ and had
to be discarded.

The fix uses local dimension 2 and then reconstructs the operator-space
quantities \emph{exactly}. For $\rho=|\psi\rangle\langle\psi|$ the operator
Schmidt values across a cut are the pairwise products of the state's Schmidt
values $\{p_i\}$, so
\begin{equation}
  \{\lambda^{\mathrm{op}}\} = \{p_ip_j\},\qquad
  \Sop = 2S_{\mathrm{vN}},\qquad
  \chreq^{\mathrm{op}}(\epsilon)\ \text{from sorting}\ \{p_ip_j\}.
  \label{eq:pureop}
\end{equation}
This gives the bond dimension an MPDO code \emph{would} have needed, from a
calculation that never pays MPDO cost, and it makes purity exact by
construction. Energy conservation via an MPO then provides a sharp truncation
diagnostic.

%======================================================================
\section{Metrics, and why we use two}
\label{sec:metrics}

\paragraph{Operator entanglement $\Sop$.} The von~Neumann entropy of the
normalised squared Schmidt spectrum of $|\rho\rangle\rangle$ across a cut,
in \emph{bits}, so that $\Sop=\log_2\chi$ is the exact statement of ``a bond
dimension $\chi$ can support this much''. We report the middle cut and the
maximum over cuts; measured profiles show these agree to within 4\%, so the
middle cut is the bottleneck.

\paragraph{Required bond dimension $\chreq(\epsilon)$.} The smallest number of
Schmidt values whose discarded weight falls below $\epsilon$ (we use $10^{-6}$
and $10^{-10}$). This, not the entropy, sets the runtime and memory.

\paragraph{Why both.} They are not interchangeable, and the recent literature
turns on the difference. Ref.~\cite{vilchez2025direct} reports that the operator
entanglement can \emph{fall} with increasing $\gamma$ while the bond dimension
actually required does not, because the entropy characterises the final state
whereas the simulation must carry the whole trajectory. Our data confirms the
gap quantitatively: Fig.~\ref{fig:growth}(c) shows
\begin{equation}
  \boxed{\ \chreq \;\approx\; 6\cdot 2^{\,1.61\,\Sop}\ }
  \label{eq:chiS}
\end{equation}
over three decades in $\chi$ with residual RMS $0.03$ bits. The naive $2^{\Sop}$
underestimates the cost by a factor $2^{0.61\Sop}$, which is $\times64$ at
$\Sop=10$ bits. \emph{Do not size a calculation from the entropy alone.}

\paragraph{Validity diagnostics.} Every recorded point carries
$\mathrm{Tr}\rho$ (must be 1), $\mathrm{Tr}\rho^2$ (must lie in $[0,1]$ and
$\to1$ as the state relaxes to $|0\cdots0\rangle$), and a \texttt{saturated}
flag that fires whenever $\chi_{\max}$ binds at any bond. The flag matters
because $\chreq$ is read off the \emph{stored} spectrum: if truncation is
active, $\chreq$ cannot exceed $\chi_{\max}$ and the measurement is
\textbf{censored}, not converged. In an early run we measured
$\chreq = 64,\,127,\,249$ at $\chi_{\max}=64,\,128,\,256$ --- tracking the cap,
not converging. Any $\chreq$ from a saturated row is a lower bound only.

%======================================================================
\section{What we expected, from the literature}
\label{sec:lit}

\paragraph{Closed system.} Ref.~\cite{prosen2007integrability} is the canonical
statement: the MPO rank needed to represent time-evolved observables grows as
$\exp(\mathrm{const}\cdot t)$ for non-integrable chains and only polynomially
when the model is integrable. Our disordered-bond chain is non-integrable; bond
disorder alone does not localise (that would need random fields).

\paragraph{Noise as a simulability threshold.} Ref.~\cite{noh2020efficient}
established that above a characteristic system size the maximum achievable MPO
entanglement is bounded by a constant depending only on the error rate, not on
system size, though that bound grows rapidly as the rate decreases.
Ref.~\cite{azad2023phase} finds an actual transition as environment coupling
increases, beyond which classical simulation succeeds at all times.

\paragraph{Closest prior work.} Ref.~\cite{preisser2023comparing} simulates an
XXZ chain under spontaneous emission and absorption from a N\'eel state and
compares operator entanglement (MPDO) against trajectory entanglement. It is
essentially our model. Their reported behaviour: fast growth and inefficiency at
$\gamma/J=0.01$; a clear rise-and-fall at $\gamma/J\gtrsim0.1$; and, for the
maximally imbalanced case (emission only, i.e.\ ours), relaxation to a trivial
product steady state so that $\Sop\to0$.

\paragraph{The prediction we carried in, and got wrong.} From the rise-and-fall
picture we assumed the barrier peaks at $t\sim1/\gamma$ and that the hard regime
was therefore \emph{small $\gamma$ at late times}. Both halves of that turned
out to be wrong; see Sec.~\ref{sec:wrong}. The cost of the error was concrete:
the first production sweep set $t_{\max}\sim2/\gamma$, which at $\gamma=0.01$
meant integrating to $t=200$ to capture a peak at $t=4$, and the jobs died on
walltime before reaching large $n$.

%======================================================================
\section{Numerical results}

All runs below: uniform $J=0.5$, N\'eel initial state, \texttt{:strang}
splitting, $dt=0.05$, cutoff $10^{-12}$, $\chi_{\max}=256$ for the open runs.

\subsection{The barrier is real, and it is transient}

\begin{figure}[htbp]
\centering
\includegraphics[width=\textwidth]{figs/fig1_barrier.pdf}
\caption{Operator entanglement versus time for each dissipation rate. The
characteristic shape is a \emph{barrier}: rise, peak, decay to zero, as required
by the trivial product steady state of amplitude damping. Crosses mark points
where $\mathrm{Tr}\rho<0.95$, i.e.\ where $\chi_{\max}=256$ was too small and the
data is untrustworthy. Note the $\gamma=0.10$ and $\gamma=0.20$ panels, where
curves for different $n$ lie on top of one another.}
\label{fig:barrier}
\end{figure}

Fig.~\ref{fig:barrier} confirms the barrier at every $\gamma$. The measured peak
times are $t^\ast \approx 4.4$--$4.8$ ($\gamma=0.10$), $4.8$--$8.0$
($\gamma=0.05$), against $1/\gamma = 10$ and $20$. The peak location is set by
the \emph{system size}, $t^\ast\approx0.45\,n$, not by the dissipation rate,
except at large $\gamma$ where dissipation cuts the growth off first.

\subsection{Dissipation is invisible until $\gamma t\approx0.4$}

\begin{figure}[htbp]
\centering
\includegraphics[width=0.86\textwidth]{figs/fig4_method.pdf}
\caption{(a) Ratio of open to closed operator entanglement at $n=8$, plotted
against $\gamma t$. Four dissipation rates spanning two decades collapse onto a
single curve which stays at unity until $\gamma t\approx0.4$ and then falls off
a cliff. (b) Measured effective Trotter order (Sec.~\ref{sec:trotter}); the
missing \texttt{:project} point at $t=10$ is where its error changes sign.}
\label{fig:method}
\end{figure}

Fig.~\ref{fig:method}(a) is the cleanest scaling result in the study: the ratio
$\Sop^{\mathrm{open}}/\Sop^{\mathrm{closed}}$ collapses across
$\gamma\in[0.01,0.2]$ and remains indistinguishable from unity until
$\gamma t\approx0.4$. Before that point the amplitude damping is doing nothing
measurable --- the simulation is a closed-system quench.

Combining $t_{\mathrm{dep}}\approx0.4/\gamma$ with the closed-system saturation
time $t_{\mathrm{sat}}\approx n/2$ gives the condition for dissipation, rather
than finite size, to cap the entanglement:
\begin{equation}
  \frac{0.4}{\gamma}\lesssim\frac{n}{2}
  \qquad\Longleftrightarrow\qquad
  \gamma\gtrsim\frac{0.8}{n}.
\end{equation}

\subsection{The area law, and the saturation size}

\begin{figure}[htbp]
\centering
\includegraphics[width=\textwidth]{figs/fig2_saturation.pdf}
\caption{(a) Peak operator entanglement versus $n$. At $\gamma=0.20$ and
$\gamma=0.10$ the peak stops moving; at $\gamma\le0.05$ it is still climbing at
the largest $n$ reached. (b) The same data as increments: zero means adding
qubits is classically free. (c) Operating diagram. The line
$\nsat=1.4/\gamma$ is calibrated on the two saturating rates; stars mark the
recommended operating points.}
\label{fig:saturation}
\end{figure}

\begin{table}[htbp]
\centering
\caption{Peak $\Sop$ [bits] versus $n$ and $\gamma$, with the minimum trace over
the run as a validity indicator. Entries with $\mathrm{Tr}\rho<0.95$ (bold) are
truncation-corrupted and are lower bounds at best. $n=8$ is exact
(MPDO ceiling $4^4=256$ equals $\chi_{\max}$).}
\label{tab:peaks}
\begin{tabular}{lcccccc}
\toprule
$\gamma$ & $n=8$ & $n=12$ & $n=16$ & $n=20$ & $n=24$ & $\min\mathrm{Tr}\rho$ (worst $n$)\\
\midrule
0.02 & 4.50 & 5.51 & \textbf{6.36} & \textbf{6.89} & --- & 0.644 \\
0.05 & 4.48 & 5.51 & 6.26 & \textbf{6.70} & --- & 0.849 \\
0.10 & 3.96 & 4.86 & 4.89 & 4.90 & 4.90 & 0.978 \\
0.20 & 2.93 & 2.92 & 2.92 & 2.92 & --- & $>0.999$ \\
\bottomrule
\end{tabular}
\end{table}

Table~\ref{tab:peaks} contains the central result.

\begin{keybox}
\textbf{At $\gamma=0.10$ the peak is $4.86,4.89,4.90,4.90$ bits for
$n=12,16,20,24$.} It does not move over a doubling of the system size. The
bond-resolved profile is a flat bulk plateau with edge structure only, the trace
stays above $0.978$, and the purity correctly returns towards $1$ as the state
relaxes to $|0\cdots0\rangle$. This is a clean, verified area law: the classical
cost is independent of $n$, so a hardware experiment at $\gamma=0.10$ is
classically reproducible at \emph{any} system size.
\end{keybox}

With $\gamma=0.20$ saturating by $n=8$ and $\gamma=0.10$ by $n\approx16$, we
calibrate
\begin{equation}
  \boxed{\ \nsat(\gamma)\;\approx\;\frac{1.4}{\gamma}\ }
\end{equation}
Below this line the peak is set by finite size; above it, by $\gamma$.

A second observation from Table~\ref{tab:peaks}: $\gamma=0.02$ and $\gamma=0.05$
are essentially \emph{indistinguishable} at every $n$ we reached
($4.50/4.48$, $5.51/5.51$, $6.36/6.26$). Below the saturation line the peak
does not depend on $\gamma$ at all, so reducing $\gamma$ from $0.05$ to $0.02$
buys no extra classical difficulty while making the dissipation nearly
invisible. There is no reason to run at $\gamma=0.02$.

\subsection{Growth laws and the cost wall}

\begin{figure}[htbp]
\centering
\includegraphics[width=\textwidth]{figs/fig3_growth.pdf}
\caption{(a,b) Closed-system growth of $\Sop$ and $\chreq$. Filled markers are
the pre-saturation window, where the run is exact (below both the truncation cap
and the physical ceiling, and before the light cone reaches the boundary); faded
lines are the truncated continuation. The curves for all $n$ collapse, as they
must while the light cone is still inside the chain. (c) The relation
Eq.~\eqref{eq:chiS} between required bond dimension and entropy. (d)
Extrapolation along the operating line, against measured lower bounds.}
\label{fig:growth}
\end{figure}

In the pre-saturation window the closed-system data collapses across $n$ and
gives, with per-$n$ slopes $1.34,1.34,1.39,1.51$ and $2.14,2.14,2.24,2.45$
respectively,
\begin{equation}
  \Sop(t)\approx 1.39\,t\ \ \text{bits},
  \qquad
  \log_2\chreq(t)\approx 2.24\,t + 2.7 .
  \label{eq:growth}
\end{equation}
The $n$-independence is a consistency check rather than a result: until the
light cone reaches the boundary the chain cannot know its own length.
Eq.~\eqref{eq:chiS} follows from the ratio of the two slopes.

\begin{table}[htbp]
\centering
\caption{\textbf{Measured} (not extrapolated) closed-system requirements at
saturation. Every $\chreq$ is capped by our own $\chi_{\max}$ and is therefore a
lower bound. Memory assumes a dense MPDO, $n\cdot4\cdot\chi^2\cdot16$ bytes.}
\label{tab:measured}
\begin{tabular}{lccc}
\toprule
$n$ & $\Sop$ at saturation [bits] & $\chreq\ \ge$ & MPDO memory $\ge$ \\
\midrule
12 & 5.6 & 1\,940 & 3 GB \\
16 & 7.2 & 18\,936 & 367 GB \\
20 & 8.5 & 44\,881 & 2.6 TB \\
24 & 10.0 & 62\,256 & 6.0 TB \\
\bottomrule
\end{tabular}
\end{table}

\subsection{Extrapolation to hardware sizes}

Along the operating line $\gamma^\ast=1.4/n$, $t^\ast=0.45n$, combining
Eqs.~\eqref{eq:growth} and \eqref{eq:chiS}:

\begin{table}[htbp]
\centering
\caption{Extrapolated requirement at the barrier peak along the operating line.
Feasible today is roughly $\chi\sim10^4$--$10^5$.}
\label{tab:extrap}
\begin{tabular}{lccccc}
\toprule
$n$ & $\gamma^\ast$ & $t^\ast$ & $\Sop$ [bits] & $\chreq$ & MPDO memory \\
\midrule
16  & 0.088 & 7.2  & 10.0 & $2^{19}$  & $10^{14}$ B \\
20  & 0.070 & 9.0  & 12.5 & $2^{23}$  & $10^{17}$ B \\
24  & 0.058 & 10.8 & 15.0 & $2^{27}$  & $10^{19}$ B \\
30  & 0.047 & 13.5 & 18.8 & $2^{33}$  & $10^{23}$ B \\
50  & 0.028 & 22.5 & 31.3 & $2^{53}$  & $10^{36}$ B \\
100 & 0.014 & 45.0 & 62.6 & $2^{104}$ & $10^{66}$ B \\
\bottomrule
\end{tabular}
\end{table}

The extrapolation is fitted on $t\lesssim4$ and read out at $t\approx10$--$22$,
so it should be treated as an upper envelope; growth bends over as the light
cone reaches the boundary, which at $n=50$ happens near $t\approx25$, just past
$t^\ast=22.5$. Even halving the exponent, however, leaves $n=24$ out of reach.
The measured lower bounds in Table~\ref{tab:measured} make the same point
without any extrapolation at all.

%======================================================================
\section{Where our initial expectations were wrong}
\label{sec:wrong}

Recorded explicitly, because two of these cost real compute time.

\begin{table}[htbp]
\centering
\caption{Corrections log.}
\begin{tabular}{p{0.30\textwidth}p{0.30\textwidth}p{0.32\textwidth}}
\toprule
\textbf{Expected} & \textbf{Measured} & \textbf{Consequence} \\
\midrule
Barrier peaks at $t\sim1/\gamma$; the hard regime is late times.
& Peak at $t^\ast\approx0.45n$, set by system size. At $\gamma=0.01,n=8$: peak at $t=4$, $1/\gamma=100$.
& First sweep integrated to $t=200$ for a peak at $t=4$; $\sim$95\% of walltime wasted, jobs died before large $n$. \\
\addlinespace
Small $\gamma$ is the hard regime.
& $\gamma$ only ever \emph{reduces} hardness. $\gamma=0.02$ and $0.05$ give identical peaks below saturation.
& ``Dissipation matters'' and ``classically hard'' are in tension; the useful region is the boundary $\gamma\approx1.4/n$. \\
\addlinespace
$\chi_{\max}=256$ ample, since $\Sop$ converged at $\chi=64$.
& That was measured at $\gamma=0.2$, the easy corner. At $\gamma=0.02,n=20$ the trace falls to $0.644$.
& $\gamma\le0.05$, $n\ge16$ data is corrupt near the peak; only $\gamma=0.10$ is fully trustworthy. \\
\addlinespace
The order-2 splitting is second order.
& $p_{\mathrm{eff}}\approx1.5$, with the error changing sign (Table~\ref{tab:order}).
& Switched to a palindromic splitting; possible knock-on for MPF work. \\
\addlinespace
Closed baseline can be run in Liouville space.
& $\mathrm{Tr}\rho^2$ reached 24.2 at $n=20$.
& Rewritten as a pure-state calculation with exact operator-space reconstruction, Eq.~\eqref{eq:pureop}. \\
\bottomrule
\end{tabular}
\end{table}

%======================================================================
\section{Recommended operating points}
\label{sec:operating}

\begin{keybox}
The operating line is
\[
  \gamma^\ast \approx \frac{1.4}{n}, \qquad t^\ast \approx 0.45\,n .
\]
\textbf{$n=24$:} $\gamma\approx0.058$, $t\approx11$. Already beyond single-node
MPDO by the measured bounds of Table~\ref{tab:measured}.\\
\textbf{$n=50$:} $\gamma\approx0.028$, $t\approx22$. Extrapolated
$\chreq\sim10^{16}$.
\end{keybox}

Three points deserve emphasis for whoever configures the hardware run.

\begin{enumerate}[leftmargin=1.5em,itemsep=3pt]
\item \textbf{$\gamma$ must scale with $n$.} Running $n>50$ at $\gamma=0.05$
  would place the experiment \emph{above} the saturation line
  ($\nsat(0.05)\approx28$), in the area-law region where a classical MPDO
  handles any system size at fixed cost. Larger $n$ at fixed $\gamma$ is not
  automatically harder; past $\nsat$ it is not harder at all.
\item \textbf{$t\approx0.45n$, not the few Trotter steps used previously.}
  Earlier project runs used $t=3$, which at any of these $n$ is on the rising
  edge, far short of the peak.
\item \textbf{Our $dt=0.05$ is a convergence choice, not a hardware choice.} It
  makes the classical run a faithful integrator of Eq.~\eqref{eq:lindblad}. The
  hardware will run a Trotterised circuit at whatever step the gate budget
  allows, which is a genuinely different map with its own entanglement growth.
  If the claim is to be ``no classical method could match my circuit'', the
  classical baseline must be rerun at the hardware's Trotter step.
\end{enumerate}

%======================================================================
\section{What this does \emph{not} establish}
\label{sec:limits}

\begin{warnbox}
\textbf{This study bounds one classical method.} Everything above concerns the
vectorised MPDO/TEBD route. The standard alternative for open systems ---
stochastic unravelling into quantum trajectories, i.e.\ pure-state MPS at local
dimension 2 averaged over realisations --- is completely untested here.
Ref.~\cite{preisser2023comparing} found that trajectory entanglement and
operator entanglement behave differently, and can move in opposite directions as
parameters vary. A regime where the MPDO route needs $\chi=10^6$ may be one
where a few thousand trajectories at $\chi=200$ reproduce $\langle Z_j\rangle$ to
within shot noise. Since the hardware experiment will be judged on observables,
not on $\rho$, that is the comparison that decides the claim.
\end{warnbox}

Secondary limitations:
\begin{itemize}[leftmargin=1.4em,itemsep=2pt]
\item The growth laws Eq.~\eqref{eq:growth} are fitted on 3--4 points per $n$ in
  the window $t\lesssim4$. The per-$n$ consistency is reassuring but the fit
  deserves a dedicated large-$\chi_{\max}$ closed-system run.
\item $\nsat=1.4/\gamma$ is calibrated on two saturating rates ($0.10$ and
  $0.20$). A third, at $\gamma\approx0.07$ with $n$ up to 24, would test it.
\item $\chreq$ is not converged anywhere except $n=8$. All larger-$n$ values are
  censored lower bounds.
\item All runs use a uniform $J$. The disorder-averaged behaviour of the
  original $\mathcal{U}[1/4,3/4]$ model has not been checked.
\end{itemize}

\paragraph{Suggested next steps, in priority order.}
(1)~A trajectory-entanglement study --- the missing half, and cheap, since it
works with pure states.
(2)~A closed-system-only run at large $\chi_{\max}$ for $n=24$--$30$ to pin
Eq.~\eqref{eq:growth} out to $t\approx12$ instead of $t\approx4$.
(3)~A $\chi_{\max}$ ladder at $n=16$, $\gamma=0.05$ to quantify how much the
$\chi_{\max}=256$ truncation actually cost at the peak.
(4)~No further runs at $\gamma=0.02$.

%======================================================================
\begin{thebibliography}{9}
\small
\bibitem{preisser2023comparing}
G.~Preisser, D.~Wellnitz, T.~Botzung and J.~Schachenmayer,
\emph{Comparing bipartite entropy growth in open-system matrix-product
simulation methods}, arXiv:2303.09426; Phys.\ Rev.\ A \textbf{108}, 042605 (2023).

\bibitem{noh2020efficient}
K.~Noh, L.~Jiang and B.~Fefferman,
\emph{Efficient classical simulation of noisy random quantum circuits in one
dimension}, Quantum \textbf{4}, 318 (2020).

\bibitem{vilchez2025direct}
E.~Vilchez-Estevez, Y.~Yosifov and Z.~Sun,
\emph{Direct probing of the simulation complexity of open quantum many-body
dynamics}, arXiv:2508.19959 (2025).

\bibitem{prosen2007integrability}
T.~Prosen and M.~\v{Z}nidari\v{c},
\emph{Is the efficiency of classical simulations of quantum dynamics related to
integrability?}, Phys.\ Rev.\ E \textbf{75}, 015202(R) (2007).

\bibitem{azad2023phase}
U.~Azad, A.~Hallam, J.~Morley and A.~G.~Green,
\emph{Phase transitions in the classical simulability of open quantum systems},
Sci.\ Rep.\ \textbf{13}, 8866 (2023).

\bibitem{fefferman2024nonunital}
B.~Fefferman, S.~Ghosh, M.~Gullans, K.~Kuroiwa and K.~Sharma,
\emph{Effect of non-unital noise on random-circuit sampling},
PRX Quantum \textbf{5}, 030317 (2024).
\end{thebibliography}

\end{document}
